\section{Protocole expérimental}
\label{sec:experiments}

\subsection{Comparaison contrôlée}

L'expérience principale compare GQA dense à l'attention contextuelle W/R/V. Les deux modèles ont dix couches, une dimension cachée de 384, huit têtes de lecture/query, deux têtes partagées d'écriture/valeur ou key/value, et le même chemin feed-forward à objets lexicaux et micro-experts. GQA compte 98~179~844 paramètres et W/R/V 98~195~204 (écart de 0,016\%). Le résidu lexical d'écriture W/R/V est fixé à zéro et omis, tandis que RMSNorm par tête sur R et W reste active.

\subsection{Données, tokenizer et budget}

Les six runs utilisent le tokenizer \texttt{o200k\_base} de 200~019 entrées et FineWeb-Edu \texttt{sample-10BT}, des séquences de 2~048, BF16, un batch de 16 et les seeds 42--44. Chaque run effectue 3~052 étapes, soit exactement 100~007~936 tokens. Un flux déterministe de 5\% des documents est réservé à l'évaluation. L'estimation finale moyenne 32 batches (1~048~576 tokens) à l'étape 3~000. Le shard archivé est fixé à la révision \texttt{87f0914}, SHA-256 \texttt{b1ba7b2c...d78e871}.

Les deux conditions utilisent AdamW, un LR maximal de $4\times10^{-4}$, weight decay 0,1, le même warmup/decay, une loss exacte sur le vocabulaire par chunks, l'exécution eager et la même pile H200. La compilation est désactivée pour les deux, car le chemin à table lexicale liée produisait des gradients non finis avec la pile compilateur testée ; eager restait fini.

\subsection{Mesures et débit}

Nous séparons entropie croisée next-token et débit d'entraînement. Pour chaque run, le débit est la médiane des points finis après l'étape 500, en excluant les points sous 80k tokens/s dus aux évaluations, checkpoints ou pauses réseau. Nous rapportons la moyenne des trois médianes. L'écart apparié est W/R/V moins GQA à seed identique. Avec $n=3$, nous donnons l'écart-type échantillonnal et un intervalle $t$ bilatéral à 95\%, sans l'interpréter comme preuve de significativité.

\subsection{Preuves archivées}

Chaque run possède une configuration JSON réalisée, un CSV par étape, un checkpoint final, un commit source et un manifeste SHA-256. Le shard Parquet FineWeb-Edu exact de 2,15 Go est archivé pour réévaluation offline. Les anciens runs H100 WVR/W/R/V n'emploient pas le même protocole et ne servent que de chemin descriptif.

\subsection{Matériel et vérifications numériques}

Les runs contrôlés utilisent un NVIDIA H200 avec PyTorch SDPA autorisant FlashAttention. Les tests FP32 couvrent causalité, padding, formes groupées, RoPE, cache, équivalence de la projection R/W/V fusionnée, invariance lexical-off et équivalence full/incrémental. Le cache W/V croît avec le préfixe, comme un cache KV groupé ; W/R/V n'est présenté ni comme état fixe ni comme sans attention.